\documentclass[10pt,aspectratio=169,serif]{beamer}

\usepackage[fontset=none, scheme=plain]{ctex}
\IfFontExistsTF{Songti SC}{%
	\setCJKmainfont{Songti SC}[AutoFakeSlant=0.2]%
	\setCJKsansfont{Songti SC}[AutoFakeSlant=0.2]%
}{%
	\setCJKmainfont{FandolSong-Regular}[Extension=.otf,BoldFont=FandolSong-Bold,AutoFakeSlant=0.2]%
	\setCJKsansfont{FandolSong-Regular}[Extension=.otf,BoldFont=FandolSong-Bold]%
}
\usepackage{amsmath,amssymb,graphicx,tikz,booktabs,multicol,calligra}
\usepackage{pgfplots}
\pgfplotsset{compat=1.18}
\usepackage{tcolorbox}
\tcbuselibrary{skins}
\usepackage{XyygU}
\usepackage{caption}
\captionsetup[figure]{name=图,labelsep=period}

\usetikzlibrary{arrows.meta,positioning,calc}

\setlength{\abovedisplayskip}{0pt plus 2pt minus 2pt}
\setlength{\belowdisplayskip}{0pt plus 2pt minus 2pt}

\title[习题5.5]{习题5.5 三角恒等变换}
\subtitle{复习巩固 \& 综合运用 \& 拓广探索}
\author{}
\institute{}
\date{2026-10-16}

\begin{document}

\begin{frame}
  \titlepage
\end{frame}

% ===== 复习巩固 =====
\begin{frame}{第一部分：复习巩固（第1--12题）}
利用两角和差公式、二倍角公式、辅助角公式进行三角恒等变换。
\end{frame}

% --- 1 ---
\begin{frame}[t]{第1题：求 $\cos(\alpha-\beta)$}
\textbf{1.} 已知 $\sin\alpha=\dfrac23$，$\cos\beta=-\dfrac34$，$\alpha\in\left(\dfrac{\pi}{2},\pi\right)$，$\beta\in\left(\pi,\dfrac{3\pi}{2}\right)$，求 $\cos(\alpha-\beta)$ 的值。
\end{frame}

\begin{frame}[t]{第1题 解答}
\begin{tcolorbox}[colback=orange!10,colframe=orange,title=\textbf{解}]
$\alpha\in(\frac{\pi}{2},\pi)$，$\cos\alpha<0$：
\[
\cos\alpha=-\sqrt{1-\frac49}=-\frac{\sqrt5}{3}.
\]
$\beta\in(\pi,\frac{3\pi}{2})$，$\sin\beta<0$：
\[
\sin\beta=-\sqrt{1-\frac9{16}}=-\frac{\sqrt7}{4}.
\]
\[
\cos(\alpha-\beta)=\cos\alpha\cos\beta+\sin\alpha\sin\beta
=\!\left(-\frac{\sqrt5}{3}\right)\!\cdot\!\left(-\frac34\right)+\frac23\!\cdot\!\left(-\frac{\sqrt7}{4}\right)
=\frac{3\sqrt5-2\sqrt7}{12}.
\]
\end{tcolorbox}
\end{frame}

% --- 2 ---
\begin{frame}[t]{第2题：拆角 $\beta=(\alpha+\beta)-\alpha$}
\textbf{2.} 已知 $\alpha,\beta$ 都是锐角，$\cos\alpha=\dfrac17$，$\cos(\alpha+\beta)=-\dfrac{11}{14}$，求 $\cos\beta$ 的值。（提示：$\beta=(\alpha+\beta)-\alpha$。）
\end{frame}

\begin{frame}[t]{第2题 解答}
\begin{tcolorbox}[colback=orange!10,colframe=orange,title=\textbf{解}]
$\alpha$ 锐角，$\sin\alpha=\sqrt{1-\frac1{49}}=\frac{4\sqrt3}{7}$。
$\alpha+\beta\in(0,\pi)$，由 $\cos(\alpha+\beta)=-\frac{11}{14}<0$ 知 $\alpha+\beta$ 在第二象限：
\[
\sin(\alpha+\beta)=\sqrt{1-\frac{121}{196}}=\frac{5\sqrt3}{14}.
\]
由 $\beta=(\alpha+\beta)-\alpha$：
\[
\cos\beta=\cos[(\alpha+\beta)-\alpha]=\cos(\alpha+\beta)\cos\alpha+\sin(\alpha+\beta)\sin\alpha
=-\frac{11}{14}\cdot\frac17+\frac{5\sqrt3}{14}\cdot\frac{4\sqrt3}{7}
=\frac{-11+60}{98}=\frac{49}{98}=\frac12.
\]
\end{tcolorbox}
\end{frame}

% --- 3 ---
\begin{frame}[t]{第3题：求 $\cos\alpha$}
\textbf{3.} 已知 $\sin(30^\circ+\alpha)=\dfrac35$，$60^\circ<\alpha<150^\circ$，求 $\cos\alpha$ 的值。
\end{frame}

\begin{frame}[t]{第3题 解答}
\begin{tcolorbox}[colback=orange!10,colframe=orange,title=\textbf{解}]
由 $60^\circ<\alpha<150^\circ$，得 $90^\circ<30^\circ+\alpha<180^\circ$，故 $\cos(30^\circ+\alpha)<0$：
\[
\cos(30^\circ+\alpha)=-\sqrt{1-\frac9{25}}=-\frac45.
\]
由 $\alpha=(30^\circ+\alpha)-30^\circ$：
\[
\cos\alpha=\cos[(30^\circ+\alpha)-30^\circ]
=\cos(30^\circ+\alpha)\cos30^\circ+\sin(30^\circ+\alpha)\sin30^\circ
=-\frac45\cdot\frac{\sqrt3}{2}+\frac35\cdot\frac12=\frac{3-4\sqrt3}{10}.
\]
\end{tcolorbox}
\end{frame}

% --- 4 ---
\begin{frame}[t]{第4题：$\triangle ABC$ 中求 $\cos C$}
\textbf{4.} 在 $\triangle ABC$ 中，$\sin A=\dfrac5{13}$，$\cos B=\dfrac35$，求 $\cos C$ 的值。
\end{frame}

\begin{frame}[t]{第4题 解答}
\begin{tcolorbox}[colback=orange!10,colframe=orange,title=\textbf{解}]
$\cos B=\frac35>0$，$B$ 为锐角，$\sin B=\frac45$。
$\sin A=\frac5{13}<\sin B=\frac{12}{15}=\frac45$，由大边对大角知 $A<B$，故 $A$ 为锐角：
\[
\cos A=\frac{12}{13}.
\]
由 $C=\pi-(A+B)$：
\[
\cos C=-\cos(A+B)=-(\cos A\cos B-\sin A\sin B)
=-\!\left(\frac{12}{13}\cdot\frac35-\frac5{13}\cdot\frac45\right)
=-\frac{36-20}{65}=-\frac{16}{65}.
\]
\end{tcolorbox}
\end{frame}

% --- 5 ---
\begin{frame}[t]{第5题：求 $\tan 2\alpha,\tan 2\beta$}
\textbf{5.} 已知 $\tan(\alpha+\beta)=3$，$\tan(\alpha-\beta)=5$，求 $\tan 2\alpha$，$\tan 2\beta$ 的值。
\end{frame}

\begin{frame}[t]{第5题 解答}
\begin{tcolorbox}[colback=orange!10,colframe=orange,title=\textbf{解}]
注意 $2\alpha=(\alpha+\beta)+(\alpha-\beta)$，$2\beta=(\alpha+\beta)-(\alpha-\beta)$。
\[
\tan 2\alpha=\frac{\tan(\alpha+\beta)+\tan(\alpha-\beta)}{1-\tan(\alpha+\beta)\tan(\alpha-\beta)}
=\frac{3+5}{1-15}=-\frac47.
\]
\[
\tan 2\beta=\frac{\tan(\alpha+\beta)-\tan(\alpha-\beta)}{1+\tan(\alpha+\beta)\tan(\alpha-\beta)}
=\frac{3-5}{1+15}=-\frac18.
\]
\end{tcolorbox}
\end{frame}

% --- 6 ---
\begin{frame}[t]{第6题：化简六式}
\textbf{6.} 化简：
\begin{enumerate}
  \item $\sin347^\circ\cos148^\circ+\sin77^\circ\cos58^\circ$；
  \item $\sin164^\circ\sin224^\circ+\sin254^\circ\sin314^\circ$；
  \item $\sin(\alpha+\beta)\cos(\gamma-\beta)-\cos(\beta+\alpha)\sin(\beta-\gamma)$；
  \item $\sin(\alpha-\beta)\sin(\beta-\gamma)-\cos(\alpha-\beta)\cos(\beta-\gamma)$；
  \item $\dfrac{\tan\frac{5\pi}{4}+\tan\frac{5\pi}{12}}{1-\tan\frac{5\pi}{4}\tan\frac{5\pi}{12}}$；
  \item $\dfrac{\sin(\alpha+\beta)-2\sin\alpha\cos\beta}{2\sin\alpha\sin\beta+\cos(\alpha+\beta)}$。
\end{enumerate}
\end{frame}

\begin{frame}[t]{第6题 解答}
\begin{tcolorbox}[colback=orange!10,colframe=orange,title=\textbf{解}]
(1) $347^\circ=360^\circ-13^\circ\Rightarrow\sin347^\circ=-\sin13^\circ$；$148^\circ=180^\circ-32^\circ\Rightarrow\cos148^\circ=-\cos32^\circ$。
原式 $=\sin13^\circ\cos32^\circ+\cos13^\circ\sin32^\circ=\sin45^\circ=\dfrac{\sqrt2}{2}$。

(2) $\sin164^\circ=\sin16^\circ$，$\sin224^\circ=-\sin44^\circ$，$\sin254^\circ=-\sin74^\circ=-\cos16^\circ$，$\sin314^\circ=-\sin46^\circ=-\cos44^\circ$。
原式 $=-\sin16^\circ\sin44^\circ+\cos16^\circ\cos44^\circ=\cos60^\circ=\dfrac12$。

(3) $\sin(\beta-\gamma)=-\sin(\gamma-\beta)$，$\cos(\beta+\alpha)=\cos(\alpha+\beta)$。
原式 $=\sin(\alpha+\beta)\cos(\gamma-\beta)+\cos(\alpha+\beta)\sin(\gamma-\beta)=\sin[(\alpha+\beta)+(\gamma-\beta)]=\sin(\alpha+\gamma)$。

(4) $=-[\cos(\alpha-\beta)\cos(\beta-\gamma)-\sin(\alpha-\beta)\sin(\beta-\gamma)]=-\cos[(\alpha-\beta)+(\beta-\gamma)]=-\cos(\alpha-\gamma)$。

(5) $=\tan\!\left(\dfrac{5\pi}{4}+\dfrac{5\pi}{12}\right)=\tan\dfrac{5\pi}{3}=-\sqrt3$。

(6) 分子 $=\sin(\alpha+\beta)-2\sin\alpha\cos\beta=(\sin\alpha\cos\beta+\cos\alpha\sin\beta)-2\sin\alpha\cos\beta=\cos\alpha\sin\beta-\sin\alpha\cos\beta=\sin(\beta-\alpha)$。
分母 $=2\sin\alpha\sin\beta+\cos\alpha\cos\beta-\sin\alpha\sin\beta=\cos\alpha\cos\beta+\sin\alpha\sin\beta=\cos(\alpha-\beta)$。
原式 $=\dfrac{\sin(\beta-\alpha)}{\cos(\alpha-\beta)}=\dfrac{-\sin(\alpha-\beta)}{\cos(\alpha-\beta)}=-\tan(\alpha-\beta)$。
\end{tcolorbox}
\end{frame}

% --- 7 ---
\begin{frame}[t]{第7题：求 $\sin 2\alpha,\cos 2\alpha$}
\textbf{7.} 已知 $\sin\alpha=0.80$，$\alpha\in\left(0,\dfrac{\pi}{2}\right)$，求 $\sin 2\alpha$，$\cos 2\alpha$ 的值（精确到 $0.01$）。
\end{frame}

\begin{frame}[t]{第7题 解答}
\begin{tcolorbox}[colback=orange!10,colframe=orange,title=\textbf{解}]
$\cos\alpha=\sqrt{1-0.64}=0.60$。
\[
\sin 2\alpha=2\sin\alpha\cos\alpha=2\times0.80\times0.60=0.96.
\]
\[
\cos 2\alpha=1-2\sin^2\alpha=1-2\times0.64=-0.28.
\]
\end{tcolorbox}
\end{frame}

% --- 8 ---
\begin{frame}[t]{第8题：六个恒等式证明}
\textbf{8.} 求证：
\begin{enumerate}
  \item $(\sin 2\alpha-\cos 2\alpha)^2=1-\sin 4\alpha$；
  \item $\tan\!\left(\dfrac{x}{2}+\dfrac{\pi}{4}\right)+\tan\!\left(\dfrac{x}{2}-\dfrac{\pi}{4}\right)=2\tan x$；
  \item $\dfrac{1+\sin 2\varphi}{\cos\varphi+\sin\varphi}=\cos\varphi+\sin\varphi$；
  \item $\dfrac{1-2\sin\alpha\cos\alpha}{\cos^2\alpha-\sin^2\alpha}=\dfrac{1-\tan\alpha}{1+\tan\alpha}$；
  \item $\dfrac{1-\cos 2\theta}{1+\cos 2\theta}=\tan^2\theta$；
  \item $\dfrac{1+\sin 2\theta-\cos 2\theta}{1+\sin 2\theta+\cos 2\theta}=\tan\theta$。
\end{enumerate}
\end{frame}

\begin{frame}[t]{第8题 解答}
\begin{tcolorbox}[colback=orange!10,colframe=orange,title=\textbf{证明}]
(1) 左边 $=\sin^2 2\alpha+\cos^2 2\alpha-2\sin 2\alpha\cos 2\alpha=1-\sin 4\alpha$。$\checkmark$

(2) 令 $u=\tan\dfrac{x}{2}$，则 $\tan(\frac{x}{2}\pm\frac{\pi}{4})=\dfrac{u\pm1}{1\mp u}$，相加得 $\dfrac{u+1}{1-u}+\dfrac{u-1}{1+u}=\dfrac{(u+1)^2-(1-u)^2}{1-u^2}=\dfrac{4u}{1-u^2}=2\cdot\dfrac{2u}{1-u^2}=2\tan x$。$\checkmark$

(3) $1+\sin 2\varphi=\sin^2\varphi+\cos^2\varphi+2\sin\varphi\cos\varphi=(\cos\varphi+\sin\varphi)^2$，除以 $\cos\varphi+\sin\varphi$ 即得。$\checkmark$

(4) 左边 $=\dfrac{(\cos\alpha-\sin\alpha)^2}{(\cos\alpha-\sin\alpha)(\cos\alpha+\sin\alpha)}=\dfrac{\cos\alpha-\sin\alpha}{\cos\alpha+\sin\alpha}=\dfrac{1-\tan\alpha}{1+\tan\alpha}$。$\checkmark$

(5) 左边 $=\dfrac{2\sin^2\theta}{2\cos^2\theta}=\tan^2\theta$。$\checkmark$

(6) 分子 $=2\sin^2\theta+2\sin\theta\cos\theta=2\sin\theta(\sin\theta+\cos\theta)$；分母 $=2\cos^2\theta+2\sin\theta\cos\theta=2\cos\theta(\sin\theta+\cos\theta)$。约分得 $\tan\theta$。$\checkmark$
\end{tcolorbox}
\end{frame}

% --- 9 ---
\begin{frame}[t]{第9题：由 $\sin(\alpha\pm\beta)$ 求证}
\textbf{9.} 已知 $\sin(\alpha+\beta)=\dfrac12$，$\sin(\alpha-\beta)=\dfrac13$，求证：
\begin{enumerate}
  \item $\sin\alpha\cos\beta=5\cos\alpha\sin\beta$；
  \item $\tan\alpha=5\tan\beta$。
\end{enumerate}
\end{frame}

\begin{frame}[t]{第9题 解答}
\begin{tcolorbox}[colback=orange!10,colframe=orange,title=\textbf{证明}]
由已知展开：
\[
\sin\alpha\cos\beta+\cos\alpha\sin\beta=\frac12,\qquad
\sin\alpha\cos\beta-\cos\alpha\sin\beta=\frac13.
\]
两式相加：$2\sin\alpha\cos\beta=\dfrac56\Rightarrow\sin\alpha\cos\beta=\dfrac5{12}$。
两式相减：$2\cos\alpha\sin\beta=\dfrac16\Rightarrow\cos\alpha\sin\beta=\dfrac1{12}$。

(1) $\sin\alpha\cos\beta=\dfrac5{12}=5\cdot\dfrac1{12}=5\cos\alpha\sin\beta$。$\checkmark$

(2) 两边同除以 $\cos\alpha\cos\beta$，得 $\tan\alpha=5\tan\beta$。$\checkmark$
\end{tcolorbox}
\end{frame}

% --- 10 ---
\begin{frame}[t]{第10题：求证 $\tan 2\theta=-4\tan(\theta+\frac{\pi}{4})$}
\textbf{10.} 已知 $\dfrac{1-\tan\theta}{2+\tan\theta}=1$，求证 $\tan 2\theta=-4\tan\!\left(\theta+\dfrac{\pi}{4}\right)$。
\end{frame}

\begin{frame}[t]{第10题 解答}
\begin{tcolorbox}[colback=orange!10,colframe=orange,title=\textbf{证明}]
由 $1-\tan\theta=2+\tan\theta$，得 $-2\tan\theta=1$，$\tan\theta=-\dfrac12$。
\[
\tan 2\theta=\frac{2\tan\theta}{1-\tan^2\theta}=\frac{-1}{1-\frac14}=-\frac43.
\]
\[
\tan\!\left(\theta+\frac{\pi}{4}\right)=\frac{\tan\theta+1}{1-\tan\theta}=\frac{-\frac12+1}{1+\frac12}=\frac13.
\]
故 $-4\tan(\theta+\frac{\pi}{4})=-\dfrac43=\tan 2\theta$。$\checkmark$
\end{tcolorbox}
\end{frame}

% --- 11 ---
\begin{frame}[t]{第11题：圆弧圆周角}
\textbf{11.} 已知一段圆弧所对的圆心角的正弦值等于 $\dfrac35$，求这段圆弧所对的圆周角的正弦、余弦和正切。
\end{frame}

\begin{frame}[t]{第11题 解答}
\begin{tcolorbox}[colback=orange!10,colframe=orange,title=\textbf{解}]
设圆心角为 $\alpha$，圆周角为 $\beta$，则 $\beta=\dfrac{\alpha}{2}$。$\sin\alpha=\dfrac35$。

若 $\alpha$ 为锐角，$\cos\alpha=\dfrac45$；若 $\alpha$ 为钝角，$\cos\alpha=-\dfrac45$。
\[
\sin\beta=\sin\frac{\alpha}{2},\quad \cos\beta=\cos\frac{\alpha}{2}.
\]
因圆周角必为锐角（所对圆弧为劣弧时），取半角公式正根：
\[
\sin\beta=\sqrt{\frac{1-\cos\alpha}{2}},\qquad \cos\beta=\sqrt{\frac{1+\cos\alpha}{2}}.
\]
\alert{情形一：$\cos\alpha=\frac45$}，$\sin\beta=\dfrac{\sqrt{10}}{10}$，$\cos\beta=\dfrac{3\sqrt{10}}{10}$，$\tan\beta=\dfrac13$。

\alert{情形二：$\cos\alpha=-\frac45$}，$\sin\beta=\dfrac{3\sqrt{10}}{10}$，$\cos\beta=\dfrac{\sqrt{10}}{10}$，$\tan\beta=3$。
\end{tcolorbox}
\end{frame}

% --- 12 ---
\begin{frame}[t]{第12题：辅助角化简}
\textbf{12.} 化简：
\begin{multicols}{2}
\begin{enumerate}
  \item $3\sqrt{15}\sin x+3\sqrt5\cos x$；
  \item $\dfrac32\cos x-\dfrac{\sqrt3}{2}\sin x$；
  \item $\sqrt3\sin\dfrac{x}{2}+\cos\dfrac{x}{2}$；
  \item $\dfrac{\sqrt2}{4}\sin\!\left(\dfrac{\pi}{4}-x\right)+\dfrac{\sqrt6}{4}\cos\!\left(\dfrac{\pi}{4}-x\right)$。
\end{enumerate}
\end{multicols}
\end{frame}

\begin{frame}[t]{第12题 解答}
\begin{tcolorbox}[colback=orange!10,colframe=orange,title=\textbf{解}]
(1) $A=\sqrt{135+45}=\sqrt{180}=6\sqrt5$。原式 $=6\sqrt5\sin(x+\varphi)$，其中 $\cos\varphi=\frac{3\sqrt{15}}{6\sqrt5}=\frac{\sqrt3}{2}$，$\varphi=\frac{\pi}{6}$。
原式 $=6\sqrt5\sin\!\left(x+\dfrac{\pi}{6}\right)$。

(2) $A=\sqrt{\frac94+\frac34}=\sqrt3$。原式 $=\sqrt3\cos\!\left(x+\frac{\pi}{6}\right)$。

(3) $A=\sqrt{3+1}=2$。原式 $=2\sin\!\left(\frac{x}{2}+\frac{\pi}{6}\right)$。

(4) $A=\sqrt{\frac{2}{16}+\frac{6}{16}}=\frac{\sqrt2}{2}$。原式 $=\dfrac{\sqrt2}{2}\sin\!\left(\dfrac{\pi}{4}-x+\dfrac{\pi}{3}\right)=\dfrac{\sqrt2}{2}\sin\!\left(\dfrac{7\pi}{12}-x\right)$。
\end{tcolorbox}
\end{frame}

% ===== 综合运用 =====
\begin{frame}{第二部分：综合运用（第13--17题）}
\end{frame}

% --- 13 ---
\begin{frame}[t]{第13题：$\triangle ABC$ 中求 $\angle C$}
\textbf{13.} 在 $\triangle ABC$ 中，已知 $\tan A$，$\tan B$ 是 $x$ 的方程 $x^2+p(x+1)+1=0$ 的两个实根，求 $\angle C$。
\end{frame}

\begin{frame}[t]{第13题 解答}
\begin{tcolorbox}[colback=orange!10,colframe=orange,title=\textbf{解}]
方程 $x^2+px+(p+1)=0$。由韦达定理：
\[
\tan A+\tan B=-p,\qquad \tan A\tan B=p+1.
\]
于是
\[
\tan(A+B)=\frac{\tan A+\tan B}{1-\tan A\tan B}=\frac{-p}{1-(p+1)}=\frac{-p}{-p}=1.
\]
在 $\triangle ABC$ 中 $A+B=\pi-C$，故
\[
\tan C=-\tan(A+B)=-1.
\]
$C\in(0,\pi)$，故 $\alert{C=\dfrac{3\pi}{4}}$。
\end{tcolorbox}
\end{frame}

% --- 14 ---
\begin{frame}[t]{第14题：选择题}
\textbf{14.} 在 $\triangle ABC$ 中，$B=\dfrac{\pi}{4}$，$BC$ 边上的高等于 $\dfrac13 BC$，则 $\cos A=$（\quad）。

(A) $\dfrac{3\sqrt{10}}{10}$\quad
(B) $\dfrac{\sqrt{10}}{10}$\quad
(C) $-\dfrac{\sqrt{10}}{10}$\quad
(D) $-\dfrac{3\sqrt{10}}{10}$
\end{frame}

\begin{frame}[t]{第14题 解答}
\begin{tcolorbox}[colback=orange!10,colframe=orange,title=\textbf{解}]
设 $BC=a$，$BC$ 边上的高 $AD=\dfrac{a}{3}$。
在 $\text{Rt}\triangle ABD$ 中，$\angle B=45^\circ$，故 $BD=AD=\dfrac{a}{3}$，$DC=\dfrac{2a}{3}$。
$\tan C=\dfrac{AD}{DC}=\dfrac{a/3}{2a/3}=\dfrac12$，故 $\sin C=\dfrac{1}{\sqrt5}$，$\cos C=\dfrac{2}{\sqrt5}$。
\[
\cos A=-\cos(B+C)=-(\cos B\cos C-\sin B\sin C)
=-\frac{\sqrt2}{2}\cdot\frac{2}{\sqrt5}+\frac{\sqrt2}{2}\cdot\frac{1}{\sqrt5}
=-\frac{\sqrt{10}}{10}.
\]
\alert{选 (C)}。
\end{tcolorbox}
\end{frame}

% --- 15 ---
\begin{frame}[t]{第15题：证明恒等式}
\textbf{15.} 求证：
\begin{enumerate}
  \item $3+\cos 4\alpha-4\cos 2\alpha=8\sin^4\alpha$；
  \item $\dfrac{\tan\alpha\tan 2\alpha}{\tan 2\alpha-\tan\alpha}+\sqrt3(\sin^2\alpha-\cos^2\alpha)=2\sin\!\left(2\alpha-\dfrac{\pi}{3}\right)$。
\end{enumerate}
\end{frame}

\begin{frame}[t]{第15题 解答}
\begin{tcolorbox}[colback=orange!10,colframe=orange,title=\textbf{证明}]
(1) 左边 $=3+(2\cos^2 2\alpha-1)-4\cos 2\alpha=2(\cos 2\alpha-1)^2=2(-2\sin^2\alpha)^2=8\sin^4\alpha$。$\checkmark$

(2) 注意 $\dfrac{\tan\alpha\tan 2\alpha}{\tan 2\alpha-\tan\alpha}=\dfrac{1}{\cot\alpha-\cot 2\alpha}=\dfrac{\sin\alpha\sin 2\alpha}{\sin 2\alpha\cos\alpha-\cos 2\alpha\sin\alpha}=\dfrac{\sin\alpha\cdot 2\sin\alpha\cos\alpha}{\sin\alpha}=2\sin\alpha\cos\alpha=\sin 2\alpha$。
左边 $=\sin 2\alpha-\sqrt3\cos 2\alpha=2\sin\!\left(2\alpha-\dfrac{\pi}{3}\right)$。$\checkmark$
\end{tcolorbox}
\end{frame}

% --- 16 ---
\begin{frame}[t]{第16题：是否存在锐角 $\alpha,\beta$}
\textbf{16.} 是否存在锐角 $\alpha,\beta$，使 $\alpha+2\beta=\dfrac{2\pi}{3}$，$\tan\dfrac{\alpha}{2}\tan\beta=2-\sqrt3$ 同时成立？若存在，求出 $\alpha,\beta$ 的度数；若不存在，请说明理由。
\end{frame}

\begin{frame}[t]{第16题 解答}
\begin{tcolorbox}[colback=orange!10,colframe=orange,title=\textbf{解}]
由 $\alpha+2\beta=\dfrac{2\pi}{3}$，得 $\dfrac{\alpha}{2}+\beta=\dfrac{\pi}{3}$。
令 $x=\tan\dfrac{\alpha}{2}$，$y=\tan\beta$，则 $xy=2-\sqrt3$，且
\[
\tan\!\left(\frac{\alpha}{2}+\beta\right)=\frac{x+y}{1-xy}=\sqrt3.
\]
故 $x+y=\sqrt3(1-xy)=\sqrt3(1-2+\sqrt3)=3-\sqrt3$。
$x,y$ 为方程 $t^2-(3-\sqrt3)t+(2-\sqrt3)=0$ 的两根，判别式 $\Delta=(3-\sqrt3)^2-4(2-\sqrt3)=12-6\sqrt3-8+4\sqrt3=4-2\sqrt3=(\sqrt3-1)^2$。
解得 $t_1=1,\ t_2=2-\sqrt3$。

若 $\tan\dfrac{\alpha}{2}=1\Rightarrow\dfrac{\alpha}{2}=45^\circ\Rightarrow\alpha=90^\circ$，不合锐角；
故 $\tan\dfrac{\alpha}{2}=2-\sqrt3=\tan15^\circ$，$\alpha=30^\circ$，$\beta=45^\circ$。
\alert{存在，$\alpha=30^\circ,\ \beta=45^\circ$}。
\end{tcolorbox}
\end{frame}

% --- 17 ---
\begin{frame}[t]{第17题：周期与最值}
\textbf{17.}
\begin{enumerate}
  \item 求函数 $f(x)=\sin\!\left(\dfrac{\pi}{3}+4x\right)+\sin\!\left(4x-\dfrac{\pi}{6}\right)$ 的周期和单调递增区间；
  \item 求函数 $f(x)=a\sin x+b\cos x\ (a^2+b^2\neq 0)$ 的最大值和最小值。
\end{enumerate}
\end{frame}

\begin{frame}[t]{第17题 解答}
\begin{tcolorbox}[colback=orange!10,colframe=orange,title=\textbf{解}]
(1) 注意 $\sin(4x+\frac{\pi}{3})=\cos(\frac{\pi}{6}-4x)=\cos(4x-\frac{\pi}{6})$。令 $u=4x-\frac{\pi}{6}$：
\[
f(x)=\cos u+\sin u=\sqrt2\sin\!\left(u+\frac{\pi}{4}\right)=\sqrt2\sin\!\left(4x-\frac{\pi}{12}\right).
\]
周期 $T=\dfrac{2\pi}{4}=\dfrac{\pi}{2}$。

递增区间：$2k\pi-\dfrac{\pi}{2}\le 4x-\dfrac{\pi}{12}\le 2k\pi+\dfrac{\pi}{2}$，即
\[
\frac{k\pi}{2}-\frac{5\pi}{48}\le x\le \frac{k\pi}{2}+\frac{7\pi}{48},\ k\in\mathbf{Z}.
\]

(2) $f(x)=\sqrt{a^2+b^2}\sin(x+\varphi)$，最大值 $\sqrt{a^2+b^2}$，最小值 $-\sqrt{a^2+b^2}$。
\end{tcolorbox}
\end{frame}

% ===== 拓广探索 =====
\begin{frame}{第三部分：拓广探索（第18--20题）}
\end{frame}

% --- 18 ---
\begin{frame}[t]{第18题：归纳一般规律}
\textbf{18.} 观察以下各等式：
\[
\sin^2 30^\circ+\cos^2 60^\circ+\sin30^\circ\cos60^\circ=\frac34,
\]
\[
\sin^2 20^\circ+\cos^2 50^\circ+\sin20^\circ\cos50^\circ=\frac34,
\]
\[
\sin^2 15^\circ+\cos^2 45^\circ+\sin15^\circ\cos45^\circ=\frac34.
\]
分析上述各式的共同特点，写出能反映一般规律的等式，并对等式的正确性作出证明。
\end{frame}

\begin{frame}[t]{第18题 解答}
\begin{tcolorbox}[colback=orange!10,colframe=orange,title=\textbf{解}]
观察发现：$60^\circ=30^\circ+30^\circ$，$50^\circ=20^\circ+30^\circ$，$45^\circ=15^\circ+30^\circ$，即后一角比前一角大 $30^\circ=\dfrac{\pi}{6}$。

\alert{猜想}：
\[
\sin^2\alpha+\cos^2\!\left(\alpha+\frac{\pi}{6}\right)+\sin\alpha\cos\!\left(\alpha+\frac{\pi}{6}\right)=\frac34.
\]

\textbf{证明}：令 $\beta=\alpha+\frac{\pi}{6}$，原式 $=\sin^2\alpha+\cos^2\beta+\sin\alpha\cos\beta$。
由 $\beta-\alpha=\frac{\pi}{6}$，$\cos(\beta-\alpha)=\frac{\sqrt3}{2}$。
积化和差：$2\sin\alpha\cos\beta=\sin(2\alpha+\frac{\pi}{6})+\sin(-\frac{\pi}{6})=\sin(2\alpha+\frac{\pi}{6})-\frac12$。
又 $\cos^2\beta=\frac{1+\cos(2\alpha+\frac{\pi}{3})}{2}$，$\sin^2\alpha=\frac{1-\cos 2\alpha}{2}$。
相加：$\frac12+\frac12+\frac12[\cos(2\alpha+\frac{\pi}{3})-\cos 2\alpha]+\frac12\sin(2\alpha+\frac{\pi}{6})-\frac14$
$=\frac34-\sin(2\alpha+\frac{\pi}{6})\sin\frac{\pi}{6}+\frac12\sin(2\alpha+\frac{\pi}{6})=\frac34$。$\checkmark$
\end{tcolorbox}
\end{frame}

% --- 19 ---
\begin{frame}[t]{第19题：单位圆证明和差化积}
\textbf{19.} 如图，单位圆上两点 $A(\cos\alpha,\sin\alpha)$，$B(\cos\beta,\sin\beta)$，$C$ 为弧 $AB$ 中点，$M$ 为 $AB$ 中点。利用图形证明：
\[
\frac12(\sin\alpha+\sin\beta)=\sin\frac{\alpha+\beta}{2}\cos\frac{\alpha-\beta}{2},\qquad
\frac12(\cos\alpha+\cos\beta)=\cos\frac{\alpha+\beta}{2}\cos\frac{\alpha-\beta}{2}.
\]
\end{frame}

\begin{frame}[t]{第19题 解答}
\begin{tcolorbox}[colback=orange!10,colframe=orange,title=\textbf{证明}]
$M$ 是弦 $AB$ 中点，由垂径定理 $OM\perp AB$，且 $OM$ 平分 $\angle AOB$。
$C$ 点在单位圆上对应的角为 $\dfrac{\alpha+\beta}{2}$，故
\[
M=\frac{A+B}{2}=\left(\frac{\cos\alpha+\cos\beta}{2},\frac{\sin\alpha+\sin\beta}{2}\right).
\]
$|OM|=\cos\dfrac{\alpha-\beta}{2}$（弦心距）。而 $M$ 在角 $\dfrac{\alpha+\beta}{2}$ 的终边上，故其坐标也等于
\[
|OM|\cdot\left(\cos\frac{\alpha+\beta}{2},\sin\frac{\alpha+\beta}{2}\right)
=\cos\frac{\alpha-\beta}{2}\left(\cos\frac{\alpha+\beta}{2},\sin\frac{\alpha+\beta}{2}\right).
\]
比较坐标即得两等式。$\checkmark$
\end{tcolorbox}
\end{frame}

% --- 20 ---
\begin{frame}[t]{第20题：估计 $f(\alpha)=\sin^n\alpha+\cos^n\alpha$ 取值}
\textbf{20.} 设 $f(\alpha)=\sin^n\alpha+\cos^n\alpha$，$n\in\{n\mid n=2k,\ k\in\mathbf{N}_+\}$。利用三角变换，估计 $f(\alpha)$ 在 $n=2,4,6$ 时的取值情况，进而猜想 $n$ 取一般值时 $f(\alpha)$ 的取值范围。
\end{frame}

\begin{frame}[t]{第20题 解答}
\begin{tcolorbox}[colback=orange!10,colframe=orange,title=\textbf{解}]
$n=2$：$f(\alpha)=\sin^2\alpha+\cos^2\alpha=1$，恒为 $1$。

$n=4$：$f=(\sin^2\alpha+\cos^2\alpha)^2-2\sin^2\alpha\cos^2\alpha=1-\dfrac12\sin^2 2\alpha\in\left[\dfrac12,1\right]$。

$n=6$：$f=(\sin^2\alpha+\cos^2\alpha)(\sin^4\alpha-\sin^2\alpha\cos^2\alpha+\cos^4\alpha)=1-3\sin^2\alpha\cos^2\alpha=1-\dfrac34\sin^2 2\alpha\in\left[\dfrac14,1\right]$。

\alert{猜想}：当 $n=2k$ 时，
\[
f(\alpha)\in\left[2^{1-k},\ 1\right]=\left[2^{1-\frac{n}{2}},\ 1\right].
\]
（当 $\alpha=\frac{\pi}{4}$ 时取最小值 $2^{1-k}$，当 $\alpha=0$ 或 $\frac{\pi}{2}$ 时取最大值 $1$。）
\end{tcolorbox}
\end{frame}

\begin{frame}[c]
  \begin{center}
    {\Huge\calligra Thanks!}
  \end{center}
\end{frame}

\end{document}
